% !TeX root = ../main.tex
\courseunit{2}
\emergencystretch=2em

\begin{problem}{1}
Consider \(N\) firms each with the constant-returns-to-scale production function \(Y=F\left(K,AL\right)\), or (using the intensive form) \(Y=AL\,f\left(k\right)\). Assume \(f'\left(\,\cdot\,\right)>0\), \(f''\left(\,\cdot\,\right)<0\). Assume that all firms can hire labor at wage \(wA\) and rent capital at cost \(r\), and that all firms have the same value of \(A\).
\begin{enumerate}[label=(\alph*)]
    \item Consider the problem of a firm trying to produce \(Y\) units of output at minimum cost. Show that the cost-minimizing level of \(k\) is uniquely defined and is independent of \(Y\), and that all firms therefore choose the same value of \(k\).
    \item Show that the total output of the \(N\) cost-minimizing firms equals the output that a single firm with the same production function has if it uses all the labor and capital used by the \(N\) firms.
\end{enumerate}
\end{problem}

\begin{problem}{2}[\textbf{The elasticity of substitution with constant-relative-risk-aversion utility.}]
Consider an individual who lives for two periods and whose utility is given by equation (2.43). Let \(P_1\) and \(P_2\) denote the prices of consumption in the two periods, and let \(W\) denote the value of the individual's lifetime income; thus the budget constraint is \(P_1\,C_1+P_2\,C_2=W\).
\begin{enumerate}[label=(\alph*)]
    \item What are the individual's utility-maximizing choices of \(C_1\) and \(C_2\), given \(P_1\), \(P_2\), and \(W\)?
    \item The elasticity of substitution between consumption in the two periods is
    \begin{align*}
        -\dfrac{P_1/P_2}{C_1/C_2}\dfrac{\partial\left(C_1/C_2\right)}{\partial\left(P_1/P_2\right)}
        \quad \text{or} \quad
        -\dfrac{\partial\ln\left(C_1/C_2\right)}{\partial\ln\left(P_1/P_2\right)}
    \end{align*}
    Show that with the utility function (2.43), the elasticity of substitution between \(C_1\) and \(C_2\) is \(1/\theta\).
\end{enumerate}
\end{problem}

\begin{problem}{3}
\begin{enumerate}[label=(\alph*)]
    \item Suppose it is known in advance that at some time \(t_0\) the government will confiscate half of whatever wealth each household holds at that time. Does consumption change discontinuously at time \(t_0\)? If so, why (and what is the condition relating consumption immediately before \(t_0\) to consumption immediately after)? If not, why not?
    \item Suppose it is known in advance that at \(t_0\) the government will confiscate from each household an amount of wealth equal to half of the wealth of the average household at that time. Does consumption change discontinuously at time \(t_0\)? If so, why (and what is the condition relating consumption immediately before \(t_0\) to consumption immediately after)? If not, why not?
\end{enumerate}
\end{problem}

\begin{problem}{4}
Assume that the instantaneous utility function \(u\left(C\right)\) in equation (2.2) is \(\ln C\). Consider the problem of a household maximizing (2.2) subject to (2.7). Find an expression for \(C\) at each time as a function of initial wealth plus the present value of labor income, the path of \(r\left(t\right)\), and the parameters of the utility function.
\end{problem}

\begin{problem}{5}
Consider a household with utility given by (2.2)--(2.3). Assume that the real interest rate is constant, and let \(W\) denote the household's initial wealth plus the present value of its lifetime labor income (the right-hand side of [2.7]). Find the utility-maximizing path of \(C\), given \(r\), \(W\), and the parameters of the utility function.
\end{problem}

\begin{problem}{6}[\textbf{Growth, saving, and \(r-g\).}]
\textcite{piketty2014capital} argues that a fall in the growth rate of the economy is likely to lead to an increase in the difference between the real interest rate and the growth rate. This problem asks you to investigate this issue in the context of the Ramsey--Cass--Koopmans model. Specifically, consider a Ramsey--Cass--Koopmans economy that is on its balanced growth path, and suppose there is a permanent fall in \(g\).
\begin{enumerate}[label=(\alph*)]
    \item How, if at all, does this affect the \(\dot{k}=0\) curve?
    \item How, if at all, does this affect the \(\dot{c}=0\) curve?
    \item At the time of the change, does \(c\) rise, fall, or stay the same, or is it not possible to tell?
    \item At the time of the change, does \(r-g\) rise, fall, or stay the same, or is it not possible to tell?
    \item In the long run, does \(r-g\) rise, fall, or stay the same, or is it not possible to tell?
    \item Find an expression for the impact of a marginal change in \(g\) on the fraction of output that is saved on the balanced growth path. Can one tell whether this expression is positive or negative?
    \item For the case where the production function is Cobb--Douglas, \(f\left(k\right)=k^{\alpha}\), rewrite your answer to part (f) in terms of \(\rho\), \(n\), \(g\), \(\theta\), and \(\alpha\). \par \remarkaccent{Hint:} Use the fact that \(f'\left(k^*\right)=\rho+\theta\,g\).
\end{enumerate}
\end{problem}

\begin{problem}{7}
Describe how each of the following affects the \(\dot{c}=0\) and \(\dot{k}=0\) curves in Figure 2.5, and thus how they affect the balanced-growth-path values of \(c\) and \(k\):
\begin{enumerate}[label=(\alph*)]
    \item A rise in \(\theta\).
    \item A downward shift of the production function.
    \item A change in the rate of depreciation from the value of zero assumed in the text to some positive level.
\end{enumerate}
\end{problem}

\begin{problem}{8}
Derive an expression analogous to (2.40) for the case of a positive depreciation rate.
\end{problem}

\begin{problem}{9}[\textbf{A closed-form solution of the Ramsey model.}]
\problemsource{\parencite{smith2006closed-form}}
Consider the Ramsey model with Cobb--Douglas production, \(y\left(t\right)=k\left(t\right)^{\alpha}\), and with the coefficient of relative risk aversion \(\left(\theta\right)\) and capital's share \(\left(\alpha\right)\) assumed to be equal.
\begin{enumerate}[label=(\alph*)]
    \item What is \(k\) on the balanced growth path \(\left(k^*\right)\)?
    \item What is \(c\) on the balanced growth path \(\left(c^*\right)\)?
    \item Let \(z\left(t\right)\) denote the capital-output ratio,
    \begin{align*}
    k\left(t\right)/y\left(t\right)
    \end{align*}
    and \(x\left(t\right)\) denote the consumption-capital ratio,
    \begin{align*}
    c\left(t\right)/k\left(t\right)
    \end{align*}
    Find expressions for \(\dot{z}\left(t\right)\) and \(\dot{x}\left(t\right)/x\left(t\right)\) in terms of \(z\), \(x\), and the parameters of the model.
    \item Tentatively conjecture that \(x\) is constant along the saddle path. Given this conjecture:
    \begin{enumerate}[label=(\roman*)]
        \item Find the path of \(z\) given its initial value, \(z\left(0\right)\).
        \item Find the path of \(y\) given the initial value of \(k\), \(k\left(0\right)\). Is the speed of convergence to the balanced growth path, \(d\ln\left[y\left(t\right)-y^*\right]/dt\), constant as the economy moves along the saddle path?
    \end{enumerate}
    \item In the conjectured solution, are the equations of motion for \(c\) and \(k\), (2.25) and (2.26), satisfied?
\end{enumerate}
\end{problem}

\begin{problem}{10}[\textbf{Capital taxation in the Ramsey--Cass--Koopmans model.}]
Consider a Ramsey--Cass--Koopmans economy that is on its balanced growth path. Suppose that at some time, which we will call time 0, the government switches to a policy of taxing investment income at rate \(\tau\). Thus the real interest rate that households face is now given by \(r\left(t\right)=\left(1-\tau\right)f'\left(k\left(t\right)\right)\). Assume that the government returns the revenue it collects from this tax through lump-sum transfers. Finally, assume that this change in tax policy is unanticipated.
\begin{enumerate}[label=(\alph*)]
    \item How, if at all, does the tax affect the \(\dot{c}=0\) locus? The \(\dot{k}=0\) locus?
    \item How does the economy respond to the adoption of the tax at time 0? What are the dynamics after time 0?
    \item How do the values of \(c\) and \(k\) on the new balanced growth path compare with their values on the old balanced growth path?
    \item
    \problemsource{\parencite{barro-mankiw-sala-i-martin1995capital}}
    Suppose there are many economies like this one. Workers' preferences are the same in each country, but the tax rates on investment income may vary across countries. Assume that each country is on its balanced growth path.
    \begin{enumerate}[label=(\roman*)]
        \item Show that the saving rate on the balanced growth path, \(\left(y^*-c^*\right)/y^*\), is decreasing in \(\tau\).
        \item Do citizens in low-\(\tau\), high-\(k^*\), high-saving countries have any incentive to invest in low-saving countries? Why or why not?
    \end{enumerate}
    \item Does your answer to part (c) imply that a policy of \emph{subsidizing investment} (that is, making \(\tau<0\)), and raising the revenue for this subsidy through lump-sum taxes, increases welfare? Why or why not?
    \item How, if at all, do the answers to parts (a) and (b) change if the government does not rebate the revenue from the tax but instead uses it to make government purchases?
\end{enumerate}
\end{problem}

\begin{problem}{11}[\textbf{Using the phase diagram to analyze the impact of an anticipated change.}]
Consider the policy described in Problem 2.10, but suppose that instead of announcing and implementing the tax at time 0, the government announces at time 0 that at some later time, time \(t_1\), investment income will begin to be taxed at rate \(\tau\).
\begin{enumerate}[label=(\alph*)]
    \item Draw the phase diagram showing the dynamics of \(c\) and \(k\) after time \(t_1\).
    \item Can \(c\) change discontinuously at time \(t_1\)? Why or why not?
    \item Draw the phase diagram showing the dynamics of \(c\) and \(k\) before \(t_1\).
    \item In light of your answers to parts (a), (b), and (c), what must \(c\) do at time 0?
    \item Summarize your results by sketching the paths of \(c\) and \(k\) as functions of time.
\end{enumerate}
\end{problem}

\begin{problem}{12}[\textbf{Using the phase diagram to analyze the impact of unanticipated and anticipated temporary changes.}]
Analyze the following two variations on Problem 2.11:
\begin{enumerate}[label=(\alph*)]
    \item At time 0, the government announces that it will tax investment income at rate \(\tau\) from time 0 until some later date \(t_1\); thereafter investment income will again be untaxed.
    \item At time 0, the government announces that from time \(t_1\) to some later time \(t_2\), it will tax investment income at rate \(\tau\); before \(t_1\) and after \(t_2\), investment income will not be taxed.
\end{enumerate}
\end{problem}

\begin{problem}{13}[\textbf{An interesting situation in the Ramsey--Cass--Koopmans model.}]
\begin{enumerate}[label=(\alph*)]
    \item Consider the Ramsey--Cass--Koopmans model where \(k\) at time 0 (which--as always--the model takes as given) is at the golden-rule level: \(k\left(0\right)=k^{GR}\). Sketch the paths of \(c\) and \(k\).
    \item Consider the same initial situation as in part (a), but in the version of the model that includes government purchases. Assume that \(G\) is constant and equal \(\bar{G}\). Crucially, \(\bar{G}\) is strictly less than \(f\left(k^{GR}\right)-\left(n+g\right)k^{GR}\) and strictly greater than \(f\left(k^*\right)-\left(n+g\right)k^*\) (where \(k^*\) is the level of \(k\) on the balanced growth path the economy would have if \(G\) were constant and equal to 0). Sketch the paths of \(c\) and \(k\). \par \remarkaccent{Hint:} This problem is hard but can be answered by thinking things through slowly and carefully.
\end{enumerate}
\end{problem}

\begin{problem}{14}
Consider the Diamond model with logarithmic utility and Cobb--Douglas production. Describe how each of the following affects \(k_{t+1}\) as a function of \(k_t\):
\begin{enumerate}[label=(\alph*)]
    \item A rise in \(n\).
    \item A downward shift of the production function (that is, \(f\left(k\right)\) takes the form \(B\,k^{\alpha}\), and \(B\) falls).
    \item A rise in \(\alpha\).
\end{enumerate}
\end{problem}

\begin{problem}{15}[\textbf{A discrete-time version of the Solow model.}]
Suppose \(Y_t=F\left(K_t,A_t\,L_t\right)\), with \(F\left(\,\cdot\,\right)\) having constant returns to scale and the intensive form of the production function satisfying the Inada conditions. Suppose also that \(A_{t+1}=\left(1+g\right)A_t\), \(L_{t+1}=\left(1+n\right)L_t\), and \(K_{t+1}=K_t+s\,Y_t-\delta K_t\).
\begin{enumerate}[label=(\alph*)]
    \item Find an expression for \(k_{t+1}\) as a function of \(k_t\).
    \item Sketch \(k_{t+1}\) as a function of \(k_t\). Does the economy have a balanced growth path? If the initial level of \(k\) differs from the value on the balanced growth path, does the economy converge to the balanced growth path?
    \item Find an expression for consumption per unit of effective labor on the balanced growth path as a function of the balanced-growth-path value of \(k\). What is the marginal product of capital, \(f'\left(k\right)\), when \(k\) maximizes consumption per unit of effective labor on the balanced growth path?
    \item Assume that the production function is Cobb--Douglas.
    \begin{enumerate}[label=(\roman*)]
        \item What is \(k_{t+1}\) as a function of \(k_t\)?
        \item What is \(k^*\), the value of \(k\) on the balanced growth path?
        \item Along the lines of equations (2.65)--(2.67), in the text, linearize the expression in subpart (i) around \(k_t=k^*\), and find the rate of convergence of \(k\) to \(k^*\).
    \end{enumerate}
\end{enumerate}
\end{problem}

\begin{problem}{16}[\textbf{Depreciation in the Diamond model and microeconomic foundations for the Solow model.}]
Suppose that in the Diamond model capital depreciates at rate \(\delta\), so that \(r_t=f'\left(k_t\right)-\delta\).
\begin{enumerate}[label=(\alph*)]
    \item How, if at all, does this change in the model affect equation (2.60) giving \(k_{t+1}\) as a function of \(k_t\)?
    \item In the special case of logarithmic utility, Cobb--Douglas production, and \(\delta=1\), what is the equation for \(k_{t+1}\) as a function of \(k_t\)? Compare this with the analogous expression for the discrete-time version of the Solow model with \(\delta=1\) from part (a) of Problem 2.15.
\end{enumerate}
\end{problem}

\begin{problem}{17}[\textbf{Social security in the Diamond model.}]
Consider a Diamond economy where \(g\) is zero, production is Cobb--Douglas, and utility is logarithmic.
\begin{enumerate}[label=(\alph*)]
    \item \textbf{Pay-as-you-go social security.} Suppose the government taxes each young individual an amount \(T\) and uses the proceeds to pay benefits to old individuals; thus each old person receives \(\left(1+n\right)T\).
    \begin{enumerate}[label=(\roman*)]
        \item How, if at all, does this change affect equation (2.61) giving \(k_{t+1}\) as a function of \(k_t\)?
        \item How, if at all, does this change affect the balanced-growth-path value of \(k\)?
        \item If the economy is initially on a balanced growth path that is dynamically efficient, how does a marginal increase in \(T\) affect the welfare of current and future generations? What happens if the initial balanced growth path is dynamically inefficient?
    \end{enumerate}
    \item \textbf{Fully funded social security.} Suppose the government taxes each young person an amount \(T\) and uses the proceeds to purchase capital. Individuals born at \(t\) therefore receive \(\left(1+r_{t+1}\right)T\) when they are old.
    \begin{enumerate}[label=(\roman*)]
        \item How, if at all, does this change affect equation (2.61) giving \(k_{t+1}\) as a function of \(k_t\)?
        \item How, if at all, does this change affect the balanced-growth-path value of \(k\)?
    \end{enumerate}
\end{enumerate}
\end{problem}

\begin{problem}{18}[\textbf{The basic overlapping-generations model.}]
\problemsource{\parencite{samuelson1958consumption-loan,allais1947economie}}
Suppose, as in the Diamond model, that \(L_t\) two-period-lived individuals are born in period \(t\) and that \(L_t=\left(1+n\right)L_{t-1}\). For simplicity, let utility be logarithmic with no discounting: \(U_t=\ln\left(C_{1t}\right)+\ln\left(C_{2t+1}\right)\).

The production side of the economy is simpler than in the Diamond model. Each individual born at time \(t\) is endowed with \(A\) units of the economy's single good. The good can be either consumed or stored. Each unit stored yields \(x>0\) units of the good in the following period.\footnote{Note that this is the same as the Diamond economy with \(g=0\), \(F\left(K_t,AL_t\right)=AL_t+xK_t\), and \(\delta=1\). With this production function, since individuals supply 1 unit of labor when they are young, an individual born in \(t\) obtains \(A\) units of the good. And each unit saved yields \(1+r=1+\partial F\left(K,AL\right)/\partial K-\delta=1+x-1=x\) units of second-period consumption.}

Finally, assume that in the initial period, period 0, in addition to the \(L_0\) young individuals each endowed with \(A\) units of the good, there are \(\left[1/\left(1+n\right)\right]L_0\) individuals who are alive only in period 0. Each of these ``old'' individuals is endowed with some amount \(Z\) of the good; their utility is simply their consumption in the initial period, \(C_{20}\).
\begin{enumerate}[label=(\alph*)]
    \item Describe the decentralized equilibrium of this economy. \par \remarkaccent{Hint:} Given the overlapping-generations structure, will the members of any generation engage in transactions with members of another generation?
    \item Consider paths where the fraction of agents' endowments that is stored, \(f_t\), is constant over time. What is total consumption (that is, consumption of all the young plus consumption of all the old) per person on such a path as a function of \(f\)? If \(x<1+n\), what value of \(f\) satisfying \(0\le f\le1\) maximizes consumption per person? Is the decentralized equilibrium Pareto efficient in this case? If not, how can a social planner raise welfare?
\end{enumerate}
\end{problem}

\begin{problem}{19}[\textbf{Stationary monetary equilibria in the Samuelson OLG model.}]
\problemsource{\parencite{samuelson1958consumption-loan}}
Consider the setup described in Problem 2.18. Assume that \(x<1+n\). Suppose that the old individuals in period 0, in addition to being endowed with \(Z\) units of the good, are each endowed with \(M\) units of a storable, divisible commodity, which we will call money. Money is not a source of utility.
\begin{enumerate}[label=(\alph*)]
    \item Consider an individual born at \(t\). Suppose the price of the good in units of money is \(P_t\) in \(t\) and \(P_{t+1}\) in \(t+1\). Thus the individual can sell units of endowment for \(P_t\) units of money and then use that money to buy \(P_t/P_{t+1}\) units of the next generation's endowment the following period. What is the individual's behavior as a function of \(P_t/P_{t+1}\)?
    \item Show that there is an equilibrium with \(P_{t+1}=P_t/\left(1+n\right)\) for all \(t\ge0\) and no storage, and thus that the presence of ``money'' allows the economy to reach the golden-rule level of storage.
    \item Show that there are also equilibria with \(P_{t+1}=P_t/x\) for all \(t\ge0\).
    \item Finally, explain why \(P_t=\infty\) for all \(t\) (that is, money is worthless) is also an equilibrium. Explain why this is the only equilibrium if the economy ends at some date, as in Problem 2.20(b) below. \par \remarkaccent{Hint:} Reason backward from the last period.
\end{enumerate}
\end{problem}

\begin{problem}{20}[\textbf{The source of dynamic inefficiency.}]
\problemsource{\parencite{shell1971infinity}}
There are two ways in which the Diamond and Samuelson models differ from textbook models. First, markets are incomplete: because individuals cannot trade with individuals who have not been born, some possible transactions are ruled out. Second, because time goes on forever, there are an infinite number of agents. This problem asks you to investigate which of these is the source of the possibility of dynamic inefficiency. For simplicity, it focuses on the Samuelson overlapping-generations model (see the previous two problems), again with log utility and no discounting. To simplify further, it assumes \(n=0\) and \(0<x<1\).
\begin{enumerate}[label=(\alph*)]
    \item \textbf{Incomplete markets.} Suppose we eliminate incomplete markets from the model by allowing all agents to trade in a competitive market ``before'' the beginning of time. That is, a Walrasian auctioneer calls out prices \(Q_0,Q_1,Q_2,\ldots\) for the good at each date. Individuals can then make sales and purchases at these prices given their endowments and their ability to store. The budget constraint of an individual born at \(t\) is thus \(Q_t\,C_{1t}+Q_{t+1}\,C_{2t+1}=Q_t\left(A-S_t\right)+Q_{t+1}\,xS_t\), where \(S_t\) (which must satisfy \(0\le S_t\le A\)) is the amount the individual stores.
    \begin{enumerate}[label=(\roman*)]
        \item Suppose the auctioneer announces \(Q_{t+1}=Q_t/x\) for all \(t>0\). Show that in this case individuals are indifferent concerning how much to store, that there is a set of storage decisions such that markets clear at every date, and that this equilibrium is the same as the equilibrium described in part (a) of Problem 2.18.
        \item Suppose the auctioneer announces prices that fail to satisfy \(Q_{t+1}=Q_t/x\) at some date. Show that at the first date that does not satisfy this condition the market for the good cannot clear, and thus that the proposed price path cannot be an equilibrium.
    \end{enumerate}
    \item \textbf{Infinite duration.} Suppose that the economy ends at some date \(T\). That is, suppose the individuals born at \(T\) live only one period (and hence seek to maximize \(C_{1T}\)), and that thereafter no individuals are born. Show that the decentralized equilibrium is Pareto efficient.
    \item In light of these answers, is it incomplete markets or infinite duration that is the source of dynamic inefficiency?
\end{enumerate}
\end{problem}

\begin{problem}{21}[\textbf{Explosive paths in the Samuelson overlapping-generations model.}]
\problemsource{\parencite{black1974uniqueness,brock1975perfect-foresight,calvo1978indeterminacy}}
Consider the setup described in Problem 2.19. Assume that \(x\) is zero, and assume that utility is constant-relative-risk-aversion with \(\theta<1\) rather than logarithmic. Finally, assume for simplicity that \(n=0\).
\begin{enumerate}[label=(\alph*)]
    \item What is the behavior of an individual born at \(t\) as a function of \(P_t/P_{t+1}\)? Show that the amount of his or her endowment that the individual sells for money is an increasing function of \(P_t/P_{t+1}\) and approaches zero as this ratio approaches zero.
    \item Suppose \(P_0/P_1<1\). How much of the good are the individuals born in period 0 planning to buy in period 1 from the individuals born then? What must \(P_1/P_2\) be for the individuals born in period 1 to want to supply this amount?
    \item Iterating this reasoning forward, what is the qualitative behavior of \(P_t/P_{t+1}\) over time? Does this represent an equilibrium path for the economy?
    \item Can there be an equilibrium path with \(P_0/P_1>1\)?
\end{enumerate}
\end{problem}
